% Options for packages loaded elsewhere
\PassOptionsToPackage{unicode}{hyperref}
\PassOptionsToPackage{hyphens}{url}
\PassOptionsToPackage{dvipsnames,svgnames,x11names}{xcolor}
\documentclass[
  10pt,
  fontset=fandol]{ctexart}
\usepackage{xcolor}
\usepackage[margin=16mm]{geometry}
\usepackage{amsmath,amssymb}
\setcounter{secnumdepth}{-\maxdimen} % remove section numbering
\usepackage{iftex}
\ifPDFTeX
  \usepackage[T1]{fontenc}
  \usepackage[utf8]{inputenc}
  \usepackage{textcomp} % provide euro and other symbols
\else % if luatex or xetex
  \usepackage{unicode-math} % this also loads fontspec
  \defaultfontfeatures{Scale=MatchLowercase}
  \defaultfontfeatures[\rmfamily]{Ligatures=TeX,Scale=1}
\fi
\usepackage{lmodern}
\ifPDFTeX\else
  % xetex/luatex font selection
\fi
% Use upquote if available, for straight quotes in verbatim environments
\IfFileExists{upquote.sty}{\usepackage{upquote}}{}
\IfFileExists{microtype.sty}{% use microtype if available
  \usepackage[]{microtype}
  \UseMicrotypeSet[protrusion]{basicmath} % disable protrusion for tt fonts
}{}
\makeatletter
\@ifundefined{KOMAClassName}{% if non-KOMA class
  \IfFileExists{parskip.sty}{%
    \usepackage{parskip}
  }{% else
    \setlength{\parindent}{0pt}
    \setlength{\parskip}{6pt plus 2pt minus 1pt}}
}{% if KOMA class
  \KOMAoptions{parskip=half}}
\makeatother
\setlength{\emergencystretch}{3em} % prevent overfull lines
\providecommand{\tightlist}{%
  \setlength{\itemsep}{0pt}\setlength{\parskip}{0pt}}
\usepackage{amsmath,amssymb,mathtools,needspace}
\xeCJKDeclareCharClass{CJK}{"2032,"0400->"04FF,"2160->"216B,"2460->"2473,"25A0,"25C7,"25CB,"25CF}
\IfFontExistsTF{Arial Unicode MS}{\xeCJKsetup{AutoFallBack=true}\setCJKfallbackfamilyfont{\CJKrmdefault}{Arial Unicode MS}}{}
\setcounter{secnumdepth}{0}
\setlength{\emergencystretch}{2em}
\allowdisplaybreaks
\usepackage{bookmark}
\IfFileExists{xurl.sty}{\usepackage{xurl}}{} % add URL line breaks if available
\urlstyle{same}
\hypersetup{
  pdftitle={引言},
  colorlinks=true,
  linkcolor={Maroon},
  filecolor={Maroon},
  citecolor={Blue},
  urlcolor={Blue},
  pdfcreator={LaTeX via pandoc}}

\title{引言}
\author{}
\date{}

\begin{document}
\maketitle

\subsubsection{9.4.1 引言}\label{ux5f15ux8a00}

Rutishauser（在 1958 年）利用矩阵的三角分解提出了计算矩阵特征值的 LR 算法，Francis（在 1961 年、1962 年）利用矩阵的 QR 分解建立了计算矩阵特征值的 QR 方法。

QR 方法是一种变换方法，是计算一般矩阵（中小型矩阵）全部特征值问题的最有效的方法之一。目前，QR 方法主要用来计算：（1）上 Hessenberg 阵的全部特征值问题，（2）对称三对角阵的全部特征值问题。QR 方法具有收敛快、算法稳定等特点。

对于一般矩阵 \(A\in\mathbb R^{n\times n}\)（或对称阵），首先用 Householder 方法将 \(A\) 化为上 Hessenberg 阵 \(B\)（或对称三对角阵），然后再用 QR 方法计算 \(B\) 的全部特征值问题。

事实上，除了可以用 Householder 法约化矩阵外，还可考虑用如下形式的平面旋转矩阵来约化：

\[
P_{ij}=\left(\begin{array}{*{11}{c}}
1&&&&&&&&&&\\
&\ddots&&&&&&&&&\\
&&1&&&&&&&&\\
&&&c&&&&s&&&\\
&&&&1&&&&&&\\
&&&&&\ddots&&&&&\\
&&&&&&1&&&&\\
&&&-s&&&&c&&&\\
&&&&&&&&1&&\\
&&&&&&&&&\ddots&\\
&&&&&&&&&&1
\end{array}\right),
\tag{9.4.1}
\]

矩阵原图标注：\(c,-s\) 所在列为``第 \(i\) 列''，\(s,c\) 所在列为``第 \(j\) 列''；\(c,s\) 所在行为``第 \(i\) 行''，\(-s,c\) 所在行为``第 \(j\) 行''。其中 \(c=\cos\theta\)，\(s=\sin\theta\)。

下面是用这种矩阵约化矩阵的一些结论。

\textbf{引理 1}　设 \(x=(\alpha_1,\allowbreak \alpha_2,\allowbreak \cdots,\allowbreak \alpha_i,\allowbreak \cdots,\allowbreak \alpha_j,\allowbreak \cdots,\allowbreak \alpha_n)^{\mathrm T}\)，其中 \(\alpha_i,\alpha_j\) 不全为零，则可选一平面旋转矩阵 \(P_{ij}\) 使 \(P_{ij}x=y\equiv(\alpha_1,\allowbreak \alpha_2,\allowbreak \cdots,\allowbreak \alpha_i^{(1)},\allowbreak \cdots,\allowbreak \alpha_j^{(1)},\allowbreak \cdots,\allowbreak \alpha_n)^{\mathrm T}\)，其中

\[
\alpha_i^{(1)}=\sqrt{\alpha_i^2+\alpha_j^2},
\tag{9.4.2}
\]

\[
\alpha_j^{(1)}=0,
\tag{9.4.3}
\]

\[
\begin{cases}
c=\alpha_i/\sqrt{\alpha_i^2+\alpha_j^2},\\
s=\alpha_j/\sqrt{\alpha_i^2+\alpha_j^2}.
\end{cases}
\tag{9.4.4}
\]

\href{../../source-images/pdf-251.jpeg}{原页}

\textbf{证明}　事实上，\(P_{ij}x\) 只改变 \(x\) 的第 \(i\) 个及第 \(j\) 个元素，且有

\[
\alpha_i^{(1)}=c\alpha_i+s\alpha_j,\qquad
\alpha_j^{(1)}=-s\alpha_i+c\alpha_i,
\]

于是可选 \(P_{ij}\)，使 \(\alpha_j^{(1)}=-s\alpha_i+c\alpha_j=0\)，即按式（9.4.4）求 \(c,s\)，且有式（9.4.2）及式（9.4.3）。

用平面旋转阵进行左变换可产生一算法，设给定 \(x=\begin{pmatrix}\alpha\\\beta\end{pmatrix}\)，计算 \(c=\cos\theta\)，\(s=\sin\theta\)，\(v=\sqrt{\alpha^2+\beta^2}\)，使 \(P_{ij}x=\begin{pmatrix}c&s\\-s&c\end{pmatrix}\begin{pmatrix}\alpha\\\beta\end{pmatrix}=\begin{pmatrix}v\\0\end{pmatrix}\)，为了防止溢出，将 \(x\) 规范化，有

\[
\eta\equiv\|x\|_\infty=\max\{|\alpha|,|\beta|\}\ne0,\qquad
x'=x/\eta=\begin{pmatrix}\alpha/\eta\\\beta/\eta\end{pmatrix},
\]

于是

\[
\begin{cases}c'=c,\ s'=s,\\v'=v/\eta.\end{cases}
\]

\textbf{算法 1}　设 \(x=\begin{pmatrix}\alpha\\\beta\end{pmatrix}\)，本算法产生 \(c,s\) 及 \(v\)。

\textbf{步 1}　计算 \(\eta=\max\{|\alpha|,|\beta|\}\)。

\textbf{步 2}　如果 \(\eta=0\)，则 \(c\leftarrow1\)，\(s\leftarrow0\)，转步 8。

\textbf{步 3}　\(\alpha'=\alpha/\eta\)。

\textbf{步 4}　\(\beta'=\beta/\eta\)。

\textbf{步 5}　\(v'=\sqrt{\alpha'^2+\beta'^2}\)。

\textbf{步 6}　\(c=\alpha'/v'\)，\(s=\beta'/v'\)。

\textbf{步 7}　\(v=\eta v'\)。

\textbf{步 8}　计算终止。

\textbf{定理 9.15}　如果 \(A\) 为非奇异矩阵，则存在正交矩阵 \(P_1,P_2,\cdots,P_{n-1}\)（即一系列平面旋转矩阵）使

\[
P_{n-1},\cdots,P_2P_1A=\begin{pmatrix}
r_{11}&r_{12}&\cdots&r_{1n}\\
&r_{22}&\cdots&r_{2n}\\
&&\ddots&\vdots\\
&&&r_{nn}
\end{pmatrix}\equiv R,
\tag{9.4.5}
\]

且 \(r_{ii}>0\quad(i=1,2,\cdots,n-1)\)。

\textbf{证明}　由于 \(A\) 的第 1 列一定存在 \(a_{j1}\ne0\)，于是，如果 \(a_{j1}\ne0\ (j=2,3,\cdots,n)\)，应用算法 1，即存在平面旋转矩阵 \(P_{12},P_{13},\cdots,P_{1n}\)，使

\[
P_{1n}\cdots P_{13}P_{12}A=\begin{pmatrix}
r_{11}&a_{12}^{(2)}&\cdots&a_{1n}^{(2)}\\
0&a_{22}^{(2)}&\cdots&a_{2n}^{(2)}\\
\vdots&\vdots&&\vdots\\
0&a_{n2}^{(2)}&\cdots&a_{nn}^{(2)}
\end{pmatrix}\equiv A^{(2)},
\]

且记 \(P_{1n}\cdots P_{12}=P_1\)。同理，如果 \(a_{2j}^{(2)}\ne0\ (j=3,\cdots,n)\)，应用算法 1，存在平面旋转矩阵

\begin{quote}
校注（agent 补充）：引理 1 证明首式原印 \(\alpha_j^{(1)}=-s\alpha_i+c\alpha_i\)，下一行则为 \(-s\alpha_i+c\alpha_j\)。按矩阵乘法，首式末项疑应为 \(c\alpha_j\)；正文保留原印。

校注（agent 补充）：定理 9.15 证明末行原印 \(a_{2j}^{(2)}\ne0\)。该步旋转消去第 2 列下方元素，疑应为 \(a_{j2}^{(2)}\ne0\)；正文保留原印。

校注（agent 补充）：算法 1 步 2 的零向量分支只指定 \(c,s\) 后即终止，没有显式赋值 \(v\)。按开头的 \(v=\sqrt{\alpha^2+\beta^2}\)，此分支应有 \(v=0\)；原算法步骤保留。
\end{quote}

\href{../../source-images/pdf-252.jpeg}{原页}

\(P_{23},\cdots,P_{2n}\)（记 \(P_2=P_{2n}\cdots P_{23}\)），使

\[
P_2P_1A=\begin{pmatrix}
r_{11}&a_{12}^{(2)}&a_{13}^{(2)}&\cdots&a_{1n}^{(2)}\\
&r_{22}&a_{23}^{(3)}&\cdots&a_{2n}^{(3)}\\
&&a_{33}^{(3)}&\cdots&a_{3n}^{(3)}\\
&&\vdots&&\vdots\\
&&a_{n3}^{(3)}&\cdots&a_{nn}^{(3)}
\end{pmatrix}.
\]

重复上述过程，最后得到：存在正交阵 \(P_1,P_2,\cdots,P_{n-1}\)，使式（9.4.5）成立。

\textbf{定理 9.16（矩阵的 QR 分解）}　如果 \(A\in\mathbb R^{n\times n}\) 为非奇异矩阵，则 \(A\) 可分解为一正交阵 \(Q\) 与上三角阵 \(R\) 的乘积，即 \(A=QR\)，且当 \(R\) 对角元素都为正数时分解唯一。

\textbf{证明}　由定理 9.12，存在正交阵 \(P_1,P_2,\cdots,P_{n-1}\)，使

\[
P_{n-1}\cdots P_2P_1A=R
\tag{9.4.6}
\]

为上三角阵，记

\[
Q^{\mathrm T}=P_{n-1}\cdots P_2P_1,
\]

于是式（9.4.6）为 \(Q^{\mathrm T}A=R\)，即 \(A=QR\)，其中 \(Q=P_1^{\mathrm T}P_2^{\mathrm T}\cdots P_{n-1}^{\mathrm T}\) 为正交阵。

现证唯一性：设有 \(A=Q_1R_1=Q_2R_2\)，其中 \(R_1,R_2\) 为上三角阵（显然为非奇异阵）且对角元素都为正数，\(Q_1,Q_2\) 为正交阵。于是

\[
Q_2^{\mathrm T}Q_1=R_2R_1^{-1},
\tag{9.4.7}
\]

由式（9.4.7）知上三角阵 \(R_2R_1^{-1}\) 为正交阵，故 \(R_2R_1^{-1}\) 为对角阵，即

\[
R_2R_1^{-1}=D=\operatorname{diag}(d_1,d_2,\cdots,d_n).
\]

因为 \(R_2R_1^{-1}\) 是正交阵，所以 \(D^2=I\)，又因 \(R_1,R_2\) 对角元素都为正数，故 \(d_i>0\ (i=1,2,\cdots,n)\)，即 \(D=I\)。于是 \(R_2=R_1\)，由式（9.4.7）得到 \(Q_2=Q_1\)。

\begin{center}\rule{0.5\linewidth}{0.5pt}\end{center}

\subsection{相关原页校注（agent 补充）}\label{ux76f8ux5173ux539fux9875ux6821ux6ce8agent-ux8865ux5145}

以下校注来自本节覆盖的原页，可能也涉及同页相邻小节；原文照录与校正说明分开保留。

\subsubsection{原书 PDF 页 252 的校注}\label{ux539fux4e66-pdf-ux9875-252-ux7684ux6821ux6ce8}

\href{../pages/pdf-252.md}{完整原页转录} · \href{../../source-images/pdf-252.jpeg}{原页图像}

校注记录：ch09b-E11。

\begin{quote}
校注（agent 补充）：定理 9.16 证明首句原印``由定理 9.12''，本段所需的平面旋转矩阵将 A 变成上三角阵的结论对应定理 9.15；原引用保留，列为疑似原书编号错误。
\end{quote}

\subsection{本节覆盖原页的校注记录}\label{ux672cux8282ux8986ux76d6ux539fux9875ux7684ux6821ux6ce8ux8bb0ux5f55}

下列记录按原页关联，含同页相邻小节的校注；计算时同时读取上文校注和完整原页。

\begin{itemize}
\tightlist
\item
  \texttt{ch09b-E11}：原书 PDF 页 252
\item
  \texttt{ch09b-E04}：原书 PDF 页 251
\item
  \texttt{ch09b-E05}：原书 PDF 页 251
\item
  \texttt{ch09b-E06}：原书 PDF 页 251
\end{itemize}

\end{document}
